\footnotesize
\par\smallskip\noindent\begin{minipage}{0.99\linewidth}\footnotesize\raggedright\textbf{Table legend.} \colorbox{red!20}{\strut R} = red/urgent; \colorbox{yellow!25}{\strut Y} = yellow/critical; \colorbox{green!18}{\strut G} = green/closed.\end{minipage}\par\vspace{0.25em}
\arrayrulecolor{black!22}
\begin{longtable}{L{0.03\linewidth} L{0.026\linewidth} L{0.087\linewidth} L{0.085\linewidth} L{0.052\linewidth} L{0.085\linewidth} L{0.095\linewidth} L{0.205\linewidth} L{0.13\linewidth} L{0.165\linewidth}}
\hline
\textbf{ID} & \textbf{Priority} & \textbf{Issue status} & \textbf{Addressing phase} & \textbf{Change version} & \textbf{Change timestamp} & \textbf{Finding} & \textbf{Static observation} & \textbf{Risk} & \textbf{Refactoring consequence}\\
\hline
\endfirsthead
\hline
\textbf{ID} & \textbf{Priority} & \textbf{Issue status} & \textbf{Addressing phase} & \textbf{Change version} & \textbf{Change timestamp} & \textbf{Finding} & \textbf{Static observation} & \textbf{Risk} & \textbf{Refactoring consequence}\\
\hline
\endhead
SCF01 & \cellcolor{red!20}R & open-addressed-by-phase & 1.1.3 & baseline-before-V1.0 & 2026-05-10T05:51:42+02:00 & Arithmetic/\allowbreak{}Notation.lean & 5095 lines, 903 declarations, 0 notation/\allowbreak{}syntax/\allowbreak{}macro declarations by static grep. & The notation module is an alias/\allowbreak{}plan-semantics monolith;\newline Part 2 must not import it as an API. & Split into true Notation/\allowbreak{}CoreSymbols, Legacy/\allowbreak{}PhaseAliases, and NeutralWrappers.\\
\hline
\addlinespace[0.24em]
SCF02 & \cellcolor{yellow!25}Y & open-addressed-by-phase & 1.1.4;\newline 1.1.5 & baseline-before-V1.0 & 2026-05-10T05:51:42+02:00 & Plan-semantic Smoke names & Many active derivation sites carry SM/\allowbreak{}R/\allowbreak{}Recheck/\allowbreak{}Plus/\allowbreak{}Full/\allowbreak{}Audit in file names, e.g. SM9h, SM9e, R3c1. & Without Plan Part 1, the code surface is not readable by domain role. & New public modules must be object/\allowbreak{}role based only;\newline old modules remain legacy sources.\\
\hline
\addlinespace[0.24em]
SCF03 & \cellcolor{yellow!25}Y & open-addressed-by-phase & 1.1.4;\newline 4.1-4.5 & baseline-before-V1.0 & 2026-05-10T05:51:42+02:00 & Operational SM9h core & SM9h leaves an ExecComplex discriminant derived from a quadratic window, but under historical names. & Mathematics is usable, but not Part-2-neutral. & First wrap Operational/\allowbreak{}QuadraticWindow and Operational/\allowbreak{}QuadraticDiscriminant.\\
\hline
\addlinespace[0.24em]
SCF04 & \cellcolor{yellow!25}Y & open-addressed-by-phase & 5.1-5.5;\newline 6.1-6.4 & baseline-before-V1.0 & 2026-05-10T05:51:42+02:00 & Orders /\allowbreak{} number theory & Discriminant, quadratic order, norm forms and Galois/\allowbreak{}Eisenstein special objects exist in bridge form. & Still no general D\_op -> D\_fund/\allowbreak{}conductor/\allowbreak{}order pipeline. & Build NumberTheory/\allowbreak{}* only after the operational profile is neutralized.\\
\hline
\addlinespace[0.24em]
SCF05 & \cellcolor{yellow!25}Y & open-addressed-by-phase & 6.5;\newline 10.4.2;\newline 10.4.3 & baseline-before-V1.0 & 2026-05-10T05:51:42+02:00 & Cayley-Dickson & The Structural/\allowbreak{}CayleyDickson path includes star, normSq, multiplication, alternative-law and norm-multiplicativity material. & Do not delete;\newline no free identification with sectors. & Register it as a structural consumer/\allowbreak{}bridge and keep it out of the arithmetic source chain until a bridge is proved.\\
\hline
\addlinespace[0.24em]
SCF06 & \cellcolor{red!20}R & open-addressed-by-phase & 2.3;\newline 11.1-11.5;\newline 17.1-17.5 & baseline-before-V1.0 & 2026-05-10T05:51:42+02:00 & Handoff & A\_to\_B exports carrier, Schur, regionNet, split, DtN, selectedLevel, and regime data. Inputs B/\allowbreak{}C/\allowbreak{}D/\allowbreak{}E\_to\_A exist as empty/\allowbreak{}small structures. & Pillar A could accidentally treat its own outputs as sources. & The access matrix enforces: A consumes only inputs;\newline outputs are forbidden-to-A.\\
\hline
\addlinespace[0.24em]
SCF07 & \cellcolor{yellow!25}Y & open-addressed-by-phase & 10.2;\newline 13.1-13.6;\newline 14.1-14.5 & baseline-before-V1.0 & 2026-05-10T05:51:42+02:00 & Finite/\allowbreak{}AQFT seeds & StateSpace, MatrixExponential, GibbsStateSeed, UnitaryEvolution, and ChannelSeed exist as A-internal seeds. & They can behave like parallel physics generators if not attached to an A-operator profile. & Mark them as consumers of BoundaryOperator/\allowbreak{}ExecMatrix;\newline no independent operator source.\\
\hline
\addlinespace[0.24em]
SCF08 & \cellcolor{green!18}G & closed-in-workflow & 1.1.1.1.0.1;\newline workflow\_policy & baseline-before-V1.0 & 2026-05-10T05:51:42+02:00 & Implementation preparation must be proof-rehearsed & Previous SM phases often exposed missing sources only during implementation. & The assistant may patch consumers, obstruct too early, or forget ExecComplex/\allowbreak{}import/\allowbreak{}no-go constraints. & Every phase record has preparation fields: dependency audit, source-derived theory objects, proof rehearsal, tactic plan, natural fallback source, and explicit no-go checklist.\\
\hline
\addlinespace[0.24em]
SCF09 & \cellcolor{green!18}G & closed-generated-view & 1.1.1.1.4;\newline 19.1 & baseline-before-V1.0 & 2026-05-10T05:51:42+02:00 & Integrated Lean module inventory & The SM9h snapshot contains 287 Lean modules, 724 internal import edges, zero missing internal imports, and an acyclic import graph according to build\_export\_script\_v1.8.py. & If module names and dependency structure are not visible in the planning tool, object/\allowbreak{}phase planning can drift away from the repository surface. & The planning package now carries a generated module catalog, dependency summaries, source-path links, and graph figures derived from the repository snapshot. These are generated views, not manual planning sources.\\
\hline
\addlinespace[0.24em]
SCF10 & \cellcolor{yellow!25}Y & open-comparison-guardrail & 8.1-8.8;\newline 9.1-9.5;\newline 12.1-12.8 & baseline-before-V1.0 & 2026-05-10T05:51:42+02:00 & TomS thread cross-check /\allowbreak{} external-reference intake & Thread pages 1-5 introduce prime-generation heuristics, imaginary-quadratic lattices, SL2Z/\allowbreak{}j, L-functions via modular forms, finite simple groups, Monster/\allowbreak{}Moonshine, Leech lattice and VOA/\allowbreak{}string-theory context. & These topics are established-reference material. They must not become CNNA generators, hidden axioms, imported physical ontology, or proof shortcuts;\newline several statements require precision corrections before formal use. & Register the material as source-audit and comparison-only objects unless a later phase derives the required CNNA object from A provenance. Add Leech/\allowbreak{}finite-simple-group/\allowbreak{}VOA/\allowbreak{}Moonshine comparison profiles after j/\allowbreak{}q-expansion, not before.\\
\hline
\addlinespace[0.24em]
SCF11 & \cellcolor{green!18}G & closed-routing-lock & 1.1.1.1.0;\newline 1.1.1.1.1-1.1.1.1.4 & baseline-before-V1.0 & 2026-05-10T05:51:42+02:00 & Phase 1.1 split refinement /\allowbreak{} Legacy.StatusGates quarantine & The next frontier combines three different work kinds: declaration classification, legacy enum quarantine, true notation extraction, and import firewall verification. & Keeping this as one implementation phase would recreate the previous monolith problem and make missing imports look like proof failures. & Split 1.1 into 1.1.1-1.1.5. Activate only 1.1.1 first;\newline move the 749 single-constructor status enums in 1.1.2, not during inventory.\\
\hline
\addlinespace[0.24em]
SCF12 & \cellcolor{green!18}G & closed-generated-view & 1.1.1.1.4 & baseline-before-V1.0 & 2026-05-10T05:51:42+02:00 & Declaration-level consumption gap & Module graph reachability is strong, but declaration-level semantic consumption is not equivalent to BuildAll/\allowbreak{}Umbrella import reachability. Current snapshot: 287 Lean modules, 724 internal import edges, 8860 declarations, 394 static reference edges, 1686 referenced symbols. & A declaration can look consumed because an umbrella imports its module, while it is only a status/\allowbreak{}audit marker or phase alias. Blind pruning or blind retention would both corrupt the derived-only path. & Insert Phase 1.1.1.1 and LEG5 before code movement: classify declarations as core-computational, proof-bridge, public-wrapper, legacy-status, legacy-policy, phase-alias, audit-only, blue-semantic-correction, dead-candidate, or final-output.\\
\hline
\addlinespace[0.24em]
SCF13 & \cellcolor{red!20}R & open-addressed-by-phase & 1.1.2 & baseline-before-V1.0 & 2026-05-10T05:51:42+02:00 & Single-constructor inductives are broader than status enums & A rough direct static scan finds about 872 single-constructor inductives in 99 modules;\newline the earlier 749 figure should be treated as a narrower status-enum subset, not the full population. & Moving every single-constructor inductive into Legacy.StatusGates would swallow policy/\allowbreak{}proof-bridge declarations and hide real provenance. & Phase 1.1.2 must consume the classification from 1.1.1.1/\allowbreak{}1.1.1.3. Constructor count alone is not a quarantine criterion.\\
\hline
\addlinespace[0.24em]
SCF14 & \cellcolor{red!20}R & active-next & 1.1.1.2 & baseline-before-V1.0 & 2026-05-10T05:51:42+02:00 & Blue objects are proof-carrying but not public-final & Object status B rows occur in Cayley-Dickson, Schur/\allowbreak{}DtN, and operational arithmetic families. These are already derived-only/\allowbreak{}proof-carrying in current code, but still semantically/\allowbreak{}API unfinished. & Treating blue objects as final because proofs exist would export old SM/\allowbreak{}AR semantics or wrong API boundaries into Phase 3+ and handoffs. & Add Phase 1.1.1.2 and LEG7. Every blue object must get a correction ledger entry: semantic defect, natural correction phase, allowed temporary consumers, forbidden public consumers, and exit-to-white condition.\\
\hline
\addlinespace[0.24em]
SCF15 & \cellcolor{yellow!25}Y & open-addressed-by-phase & 1.1.1.4 & baseline-before-V1.0 & 2026-05-10T05:51:42+02:00 & Pruning must follow API contracts, not precede them & Syntactic sinks and weakly referenced declarations include diagnostic leaves, audit leaves, intentional outputs, and proof-obligation terminals. & Deleting before a public API contract can remove audit/\allowbreak{}provenance evidence or proof bridges that are not mainline outputs but still necessary. & Add Phase 1.1.1.4 and LEG8. Deletion is postponed to a later optional pruning phase and requires classification plus local green build evidence.\\
\hline
\addlinespace[0.24em]
SCF16 & \cellcolor{green!18}G & closed-generated-view & 1.1.1.1.1 & baseline-before-V1.0 & 2026-05-10T05:51:42+02:00 & Module budget missing before implementation & The v23 plan classifies module/\allowbreak{}declaration consumption but did not yet require every implementation phase to name the exact modules to add/\allowbreak{}modify before Lean code is written. & Kleinteilige phases can degenerate into kleinteilige code: one new file per subphase, accidental wrappers, and parallel helper modules that are syntactically green but semantically fragmented. & Insert 1.1.1.1.1 and LEG10. Every implementation phase must carry a module manifest with exact module names/\allowbreak{}count, semantic placement rationale, reuse decision, and import role before code starts.\\
\hline
\addlinespace[0.24em]
SCF17 & \cellcolor{green!18}G & closed-generated-view & 1.1.1.1.2 & baseline-before-V1.0 & 2026-05-10T05:51:42+02:00 & Duplicate generator risk & The Smoke-Test Generator is intended to become the canonical generator seam for the whole theory, but later phases could accidentally reimplement request/\allowbreak{}output, extraction, ledger, or smoke-test logic under new names. & Duplicate generator paths would break derived-only provenance by creating competing sources of truth and would make future proof reuse fragile. & Insert 1.1.1.1.2 and LEG11. Generator-like code must reuse or extend the canonical generator at its natural layer;\newline parallel generators are blocked unless a genuine semantic split is recorded.\\
\hline
\addlinespace[0.24em]
SCF18 & \cellcolor{yellow!25}Y & deferred-secondary-roadmap & 19.1-19.4 & baseline-before-V1.0 & 2026-05-10T05:51:42+02:00 & Secondary derivation-graph analysis assets exist & repo\_inventory/\allowbreak{}raw\_export already contains import graph data and related forensics sufficient for future documentation coverage and derivation-geometry diagnostics. & If activated too early, graph metrics could be mistaken for CNNA-generative evidence or distract from closing operational derived objects. & Keep as secondary long-term phases 19.1-19.4;\newline use for documentation/\allowbreak{}audit only until core derived objects and public API are stable.\\
\hline
\addlinespace[0.24em]
SCF19 & \cellcolor{green!18}G & closed-policy-open-consumers & 1.1.1.1.0.2;\newline 1.2.1.2;\newline 4.5.1 & baseline-before-V1.0 & 2026-05-10T05:51:42+02:00 & ExecComplex rational ceiling & ExecComplex is represented as rational pairs Q x Q;\newline later CM/\allowbreak{}lattice objects require algebraic or analytic data. & Operational path may become a skeleton or silently import mirror/\allowbreak{}noncomputable semantics. & Adopt EXEC-2: operational code carries rational/\allowbreak{}integer invariants;\newline algebraic irrational data remains mirror-only unless a return invariant is proved.\\
\hline
\addlinespace[0.24em]
SCF20 & \cellcolor{green!18}G & closed-policy-open-consumers & 1.1.1.1.0.3;\newline 1.1.2-1.1.5 & baseline-before-V1.0 & 2026-05-10T05:51:42+02:00 & Refactor regression risk & Notation/\allowbreak{}alias cleanup touches a large declaration and import surface. & A single large move can create ambiguous red states and consumer-side patches. & Use wrapper-first bounded migrations with declaration-count, import-delta, consumer-breakage, and local full-build evidence.\\
\hline
\addlinespace[0.24em]
SCF21 & \cellcolor{green!18}G & closed-policy-open-consumer & 1.1.1.1.0.7;\newline 3.0 & baseline-before-V1.0 & 2026-05-10T05:51:42+02:00 & IdealThread singleton collapse is code-backed & ToC/\allowbreak{}Contract.lean defines IdealThread.carrier L as \{U.cell L\};\newline IdealThread.family uses that carrier;\newline deterministic coarsen\_carrier\_eq\_prev proves coarsenSet of the next carrier is the previous singleton. & Any plan that uses an existing IdealThread route as a nontrivial boundary matrix source would silently collapse to scalar/\allowbreak{}degenerate data. & Add S07 and BR06. Later full-boundary/\allowbreak{}DtN phases must explicitly reject IdealThread as full matrix source unless a separate lifting theorem is supplied.\\
\hline
\addlinespace[0.24em]
SCF22 & \cellcolor{green!18}G & closed-policy-open-consumer & 1.1.1.1.0.7;\newline 3.0 & baseline-before-V1.0 & 2026-05-10T05:51:42+02:00 & CutSpec.full is the maximal carrier route & Finite/\allowbreak{}CutSpec.lean defines full T with carrier := T.carrier and full\_carrier is rfl. & Confusing CutSpec.full with IdealThread or with a diagonal/\allowbreak{}scalar route would erase the intended sector/\allowbreak{}interior/\allowbreak{}interface geometry. & Add S08 and require Phase 3.0 boundary-route separation before neutral Schur/\allowbreak{}DtN wrappers.\\
\hline
\addlinespace[0.24em]
SCF23 & \cellcolor{green!18}G & closed-policy-open-consumers & 1.1.1.1.0.13;\newline 1.2.1.2;\newline 4.5.1;\newline 7.1-7.4 & baseline-before-V1.0 & 2026-05-10T05:51:42+02:00 & Exact quadratic scalar provenance gap & v35 treated O09 as optional/\allowbreak{}pilot;\newline CM/\allowbreak{}lattice phases will require exact quadratic coordinates, but derived-only prohibits free sqrt(D). & A scalar consumer could silently introduce theta/\allowbreak{}sqrt(D) without an operational discriminant source, turning a derived-only construction into an axiom or mirror import. & v36 requires O09 as schema, adds O10/\allowbreak{}BR14 as D-source instantiation bridge, and adds GNG19 no-free-sqrt enforcement.\\
\hline
\addlinespace[0.24em]
SCF24 & \cellcolor{green!18}G & closed-in-v44-generator & all phases before implementation;\newline immediate enforcement before 1.1.1.2 & baseline-before-V1.0 & 2026-05-10T05:51:42+02:00 & Static finding lifecycle fields were missing & The Section 2 Static code findings table listed risks and consequences but did not force an issue ID, current status, priority color, or addressing phase. & A known finding could be re-discovered in later work, silently left open, or closed by an unrelated phase without a visible provenance link. & Every static finding row now requires ID, Priority, Issue status, and Addressing phase. New findings must be added to Section 2 before implementation continues.\\
\hline
\addlinespace[0.24em]
SCF25 & \cellcolor{green!18}G & closed-in-v44-workflow-policy & all active phases;\newline current gate before 1.1.1.2 & baseline-before-V1.0 & 2026-05-10T05:51:42+02:00 & Pre-start feasibility check was implicit, not a hard visible workflow gate & The scratchpad contained dependency audit/\allowbreak{}proof rehearsal fields, but the workflow table did not require a separate go/\allowbreak{}no-go feasibility decision immediately before starting the next phase. & A phase could begin from a green prior build while a static finding or missing dependency still requires a natural backjump, producing consumer-side patches or over-fragmented governance work. & A mandatory pre-start feasibility workflow now checks active phase dependencies, static findings, module manifest, reuse audit, and whether the phase is infrastructure-only or code-changing before implementation starts.\\
\hline
\addlinespace[0.24em]
SCF26 & \cellcolor{yellow!25}Y & open-addressed-by-active-phase & 1.1.1.2 & baseline-before-V1.0 & 2026-05-10T05:51:42+02:00 & Blue correction cannot rely only on status B labels & The active Phase 1.1.1.2 consumes object register rows with Object status B, but a valid LEG7 ledger must also bind each row to proof dossier, current derivation site, allowed temporary consumers, forbidden consumers, and exit-to-white condition. & A proof-carrying blue object could be whitened or publicly wrapped merely because the object status row was found, without recording the exact semantic/\allowbreak{}API defect and correction phase. & Phase 1.1.1.2 must generate or validate a per-blue-object correction ledger, not a status-only summary. Public API and pruning phases remain blocked until that ledger is closed.\\
\hline
\addlinespace[0.24em]
SCF27 & \cellcolor{green!18}G & closed-in-tool-v1.0-legend-policy & pre-1.1.1.2 tool V1.0 documentation update & V1.0 & 2026-05-10T05:51:42+02:00 & Generated table views needed local, section-specific legends. & Earlier rendering either lacked local legends or inserted an over-broad mixed legend. The current generator emits static/\allowbreak{}phase/\allowbreak{}object legends only where the corresponding colored cells exist. & Wrong or missing legend text makes color semantics ambiguous across static findings, object status, and phase status tables. & Use explicit legend kinds in the generator and no legend for non-status tables.\\
\hline
\addlinespace[0.24em]
SCF28 & \cellcolor{green!18}G & closed-in-tool-v1.01-versioning & tooling/\allowbreak{}versioning-policy V1.01 & V1.01 & 2026-05-10T06:17:44+02:00 & Tool contract version and data snapshot version were conflated as v40-v44 document increments. & The previous split still allowed the Lean-code data timestamp to move when only the generator/\allowbreak{}PDF display changed. & A documentation-only edit could be misread as a Lean-code data change, destroying reproducible alignment between Lean snapshots and planning artifacts. & Metadata now separates tool\_version, tool\_change\_timestamp, and data\_timestamp. tool\_version increments by 0.01 for tool/\allowbreak{}display changes;\newline data\_timestamp is restored and changes only with Lean code/\allowbreak{}data.\\
\hline
\addlinespace[0.24em]
SCF29 & \cellcolor{green!18}G & closed-in-tool-v1.0-static-finding-migration & pre-1.1.1.2 tool V1.0 documentation update & V1.0 & 2026-05-10T05:51:42+02:00 & Essential findings front matter duplicated issue tracking outside Section 2. & The v44 PDF opened with prose bullets that contained actionable workflow findings, while Section 2 was supposed to be the issue ledger. & Important problems can remain unprioritized, unassigned, or forgotten if they live outside the Static code findings table. & The separate Essential findings section was removed;\newline the actionable content is represented as SCF rows with priority, lifecycle status, and addressing assignment.\\
\hline
\addlinespace[0.24em]
SCF30 & \cellcolor{green!18}G & closed-in-tool-v1.0-layout & pre-1.1.1.2 tool V1.0 documentation update & V1.0 & 2026-05-10T05:51:42+02:00 & Complete Lean module catalog: Module column could run visually into the Group column. & Long dotted Lean module names were emitted inside raw href cells without dot-level breakpoints and with too little Module-column width. & The generated module inventory becomes unreliable as a review object if readers cannot distinguish module names from group labels. & Module and path link text is now breakable at dots;\newline the Module column is widened and the Group/\allowbreak{}Path widths are adjusted.\\
\hline
\addlinespace[0.24em]
SCF31 & \cellcolor{green!18}G & closed-in-tool-v1.0-object-links & pre-1.1.1.2 tool V1.0 documentation update & V1.0 & 2026-05-10T05:51:42+02:00 & Object table A IDs were row anchors but not clickable links. & object\_table\_id\_cell(target=True) emitted a hypertarget around plain ID text, while other object tables emitted hyperlinks to that row. & The first object table looked less navigable than the rest of the proofbook workflow and did not expose the proof-dossier jump from the canonical object list. & Object table A IDs now remain row anchors and additionally link to the corresponding proof-dossier target.\\
\hline
\addlinespace[0.24em]
SCF32 & \cellcolor{green!18}G & closed-in-tool-v1.01-legend-policy & tooling/\allowbreak{}static-finding-rendering V1.01 & V1.01 & 2026-05-10T06:17:44+02:00 & Global table legends were semantically over-broad. & The V1.0 generator inserted the same mixed object/\allowbreak{}phase/\allowbreak{}static legend before generated tables even when the table had no colored status cells or the section required a narrower legend. & A reader could interpret colors with the wrong status system, especially static findings and phase rows. & The generator now supports explicit legend kinds (static, phase, object) and emits no legend for tables without colored status/\allowbreak{}triage semantics.\\
\hline
\addlinespace[0.24em]
SCF33 & \cellcolor{green!18}G & closed-in-tool-v1.01-data-timestamp-policy & tooling/\allowbreak{}versioning-policy V1.01 & V1.01 & 2026-05-10T06:17:44+02:00 & Lean-code data timestamp advanced during a display-only update. & The V1.0 legend-policy update reported a new data timestamp even though no Lean source file was changed. & Downstream reviewers could infer a new Lean data state from a purely presentational PDF/\allowbreak{}generator change. & data\_timestamp/\allowbreak{}data\_version are reset to 2026-05-10T05:51:42+02:00;\newline tool\_change\_timestamp records the current display-policy correction separately.\\
\hline
\addlinespace[0.24em]
SCF34 & \cellcolor{green!18}G & closed-in-tool-v1.01-record-change-tracking & tooling/\allowbreak{}table-rendering V1.01 & V1.01 & 2026-05-10T06:17:44+02:00 & Changed static findings, scratchpad records, and object rows had no visible per-record version/\allowbreak{}timestamp field. & Before V1.01, table readers could see final status but not which tool/\allowbreak{}display version last changed the row. & Review history becomes ambiguous when a table row changes without visible version/\allowbreak{}time metadata. & Static findings expose Change version and Change timestamp. Scratchpad and object proof-link views render record metadata and default unchanged rows to unchanged-by-V1.01 /\allowbreak{} restored Lean data timestamp.\\
\hline
\addlinespace[0.24em]
SCF35 & \cellcolor{green!18}G & closed-in-tool-v1.02-pdf-navigation & tooling/\allowbreak{}pdf-navigation V1.02 & V1.02 & 2026-05-10T06:27:26+02:00 & PDF navigation lacked a clickable contents page and explicit outline/\allowbreak{}bookmark tree. & Before V1.02, the PDF body started directly with Latest change overview and relied mainly on reader search or section scrolling. The existing hyperref setup did not expose a document-level contents page or an explicitly configured bookmark/\allowbreak{}outline tree as a reviewed navigation feature. & Long generated planning PDFs become hard to review, especially when checking sections such as Static code findings, object tables, phase tables, and scratchpad records after tool-only changes. & The PDF preamble now configures section-level PDF bookmarks/\allowbreak{}outline support, and the document body emits a clickable table of contents immediately after the title. This is a tool/\allowbreak{}display change only and does not advance the Lean-code data timestamp.\\
\hline
\addlinespace[0.24em]
SCF36 & \cellcolor{green!18}G & closed-in-tool-v1.03-pre-implementation-phase-check & tooling/\allowbreak{}pre-implementation-phase-check-workflow & V1.03 & 2026-05-10T06:38:25+02:00 & Forward-looking semantic phase feasibility checks lacked an explicit workflow contract. & The plan had per-phase preparation fields and a pre-start feasibility rule, but no dedicated workflow for simulating several semantically connected future phases before implementation and classifying them as implementable or dependency-blocked. & Without this workflow, missing dependencies can be discovered only during implementation;\newline red findings may not automatically generate upstream addressing phases, causing late backjumps and ad hoc patching pressure. & Tool V1.03 introduces a Pre-Implementation Phase Check: inspect a semantic window as a thought experiment, mark phases Y when implementable and R when dependencies are missing, record new findings, and generate upstream addressing phases for every red result before the blocked phase can become active.\\
\hline
\addlinespace[0.24em]
SCF37 & \cellcolor{green!18}G & closed-in-tool-v1.04-implementation-scaling-plan & tooling/\allowbreak{}implementation-scaling-phase-plan V1.04 & V1.04 & 2026-05-10T06:58:41+02:00 & The tool lacked an explicit phase roadmap for scalable CNNA object-world growth without overloading Lean implementation work. & Previous workflows defined phase checks and ledgers, but did not yet separate lean-focus packets, local module-theory context, large end-theory comparison ledgers, incremental cache, object graph slices, and automatic red-blocker prephases into implementation phases. & Without this roadmap, future requests could drift into manual PDF/\allowbreak{}table maintenance, excessive context scanning, or ambiguous treatment of established theories during Lean-code work. & Tool V1.04 adds future phases under 1.1.1.5 and a generated implementation-scaling plan. These phases preserve the active 1.1.1.2 cursor and keep Lean-data timestamp fixed.\\
\hline
\addlinespace[0.24em]
SCF38 & \cellcolor{green!18}G & closed-in-tool-v1.05-preimplementation-window-result-ledger & tooling/\allowbreak{}pre-implementation-check-result-ledger V1.05 & V1.05 & 2026-05-10T07:12:27+02:00 & Requested semantic phase-window implementability checks had no persistent result table. & Tool V1.03 defined the pre-implementation thought-experiment workflow, but the concrete result of a requested check window was not yet rendered as a dedicated generated ledger with phase-level Y/\allowbreak{}R decisions. & Feasibility conclusions could remain only in chat prose, making later phase activation ambiguous and weakening the red-blocker-to-addressing-phase discipline. & Tool V1.05 adds a pre\_implementation\_phase\_check\_results section, generator validation, TSV export, and PDF rendering for the checked window. The 1.1.1.2-1.1.1.5 window is recorded as yellow/\allowbreak{}implementable under sequential closure assumptions, with no red addressing phase generated.\\
\hline
\addlinespace[0.24em]
SCF39 & \cellcolor{green!18}G & closed-in-tool-v1.06-preimplementation-status-refinement & tooling/\allowbreak{}pre-implementation-check-refinement V1.06 & V1.06 & 2026-05-10T07:22:31+02:00 & Pre-implementation check status Y was too coarse and conflated active implementability with later implementability in a derivation chain. & The V1.05 check correctly treated 1.1.1.3-1.1.1.5 as implementable only under predecessor-closure assumptions, but the table still marked every implementable row yellow. & Readers could mistake later chain-dependent phases for immediately active implementable phases, weakening the one-active-phase invariant and obscuring predecessor assumptions. & Tool V1.06 introduces G/\allowbreak{}Y/\allowbreak{}O/\allowbreak{}R check semantics plus Check mode, Implementability scope, Required predecessor closures, Assumptions, Evidence basis, Lean proof risk, Parent/\allowbreak{}subphase rule, and Would-generate-addressing-phase fields.\\
\hline
\addlinespace[0.24em]
SCF40 & \cellcolor{green!18}G & closed-in-tool-v1.06-phase-legend-synchronization & tooling/\allowbreak{}phase-logic-legends V1.06 & V1.06 & 2026-05-10T07:22:31+02:00 & Tables that use phase logic needed legends synchronized with the new orange derivation-chain implementability status. & The prior phase and pre-implementation legends exposed only green/\allowbreak{}white/\allowbreak{}yellow/\allowbreak{}red or yellow/\allowbreak{}red meanings and did not explain the new chain-dependent implementability state. & Even with correct result rows, inconsistent legends would make the PDF display ambiguous and could lead to wrong activation of a non-current phase. & The generator now emits updated phase-logic legends: G = closed/\allowbreak{}derived-only, W = not yet checked where actual roadmap statuses appear, O = implementable in derivation chain, Y = active implementable now, R = blocked with missing dependencies and mandatory upstream addressing phase.\\
\hline
\addlinespace[0.24em]
SCF41 & \cellcolor{green!18}G & closed-in-tool-v1.07-db-review-main-side-objective & 20;\newline 20.1 & Tool V1.07 & 2026-05-10T07:39:10+02:00 & Database-ready review/\allowbreak{}documentation structure was only a prose idea, not a first-class main side objective. & The prior YAML had secondary long-term goals and generated views, but no main phase route for database-ready review structure, round-trip schema discipline, and expert-review semantics. & The DB direction could be treated as optional late tooling and fail to constrain current YAML/\allowbreak{}schema/\allowbreak{}review practices. & Add Phase 20 DB-REVIEW with child phases, direct DBR objects, scratchpads, secondary-goal row, workflow row, and immediate guardrails. No Lean source changes.\\
\hline
\addlinespace[0.24em]
SCF42 & \cellcolor{green!18}G & closed-in-tool-v1.07-expert-review-guardrail & 20.1;\newline GNG23;\newline workflow\_policy & Tool V1.07 & 2026-05-10T07:39:10+02:00 & Generated documentation lacked an explicit global no-infallibility and discipline-specific expert-review rule. & The tool distinguished Lean code from generated views, but did not yet force the stronger epistemic maxim that even green Lean builds and generated artifacts remain semantically reviewable. & Planning tables or green builds could be overread as 100-percent secured truth, especially around established-theory/\allowbreak{}reference material that may still lack blue marking. & Add GNG23 and the DB-REVIEW workflow: Lean code is primary but not semantically infallible;\newline generated artifacts are secondary;\newline expert review is always required;\newline missing blue markings are audit debt.\\
\hline
\addlinespace[0.24em]
SCF43 & \cellcolor{yellow!25}Y & active-in-phase-1.1.1.1.5.1 & 1.1.1.1.5.1-1.1.1.1.5.5 & V1.12 & 2026-05-10T11:16:00+02:00 & New substantive Lean code should wait until semantically related module families have been audited for consolidation or explicit no-merge status. & The repository has a valid import graph, but it still contains visibly related micro-modules and clone-like wrappers. Existing phases covered reuse/\allowbreak{}generator policy and declaration consumption, but not a dedicated pre-substantive module-consolidation/\allowbreak{}no-merge ledger. & Without this gate, new code may be built on duplicate or unstable import surfaces, making later public API correction, blue-object review, and refactor migration harder. & V1.21 routes this finding through the full readiness chain: source intake, neutral naming/\allowbreak{}facade plan, refactor execution, post-refactor re-scan, and release proof before 1.1.1.2.\\
\hline
\addlinespace[0.24em]
SCF44 & \cellcolor{yellow!25}Y & prepared-after-source-intake & 1.1.1.1.5.2-1.1.1.1.5.3 & V1.12 & 2026-05-10T11:16:00+02:00 & High-confidence schema/\allowbreak{}clone families were found in BoundarySource substrate ledgers and Handoff input/\allowbreak{}output wrappers. & BoundaryStateSubstrate, BranchAddressSubstrate, HistoryExpandedSubstrate, and InterfaceExposureSubstrate share the same import, ledger projection pattern, blocker/\allowbreak{}status/\allowbreak{}obligation theorems, and StrengtheningAR2b1 wrappers. A\_to\_C, A\_to\_D, A\_to\_E are SectorExportStrong aliases with identical of/\allowbreak{}export projections. B\_to\_A, C\_to\_A, D\_to\_A, E\_to\_A share the same source/\allowbreak{}feedback structure, ext proof, and handoff projection. & Leaving these as unrelated modules can obscure the true shared schema and make future public API/\allowbreak{}firewall decisions noisier;\newline merging them unsafely could still break existing imports or erase pillar-specific names. & After LEG37 closes, LEG38 must decide schema/\allowbreak{}facade/\allowbreak{}no-merge treatment;\newline any actual code change moves through LEG39 with old-module facades and local build evidence.\\
\hline
\addlinespace[0.24em]
SCF45 & \cellcolor{yellow!25}Y & prepared-after-source-intake & 1.1.1.1.5.2 & V1.12 & 2026-05-10T11:16:00+02:00 & Long SM/\allowbreak{}generated proof chains and Cayley-Dickson S6 chains are not safe merge candidates merely because they are numerous or similarly named. & The Smoke/\allowbreak{}GeneratedInterior*/\allowbreak{}Operational* modules and Structural/\allowbreak{}CayleyDickson/\allowbreak{}S6* modules encode phase/\allowbreak{}proof progression and consumer history. Their boundaries may carry provenance, debugging, or theorem-staging information even if some boilerplate repeats. & Aggressive consolidation could reduce auditability, hide natural proof boundaries, or create a large monolith that is harder for experts to review. & LEG38 must explicitly classify these as no-merge-until-proven unless a canonical proof module preserves provenance and passes import/\allowbreak{}declaration regression checks.\\
\hline
\addlinespace[0.24em]
SCF46 & \cellcolor{green!18}G & closed-by-V1.13-policy & 20.8 & V1.13 & 2026-05-10T11:26:00+02:00 & The project root had accumulated generated documentation artifacts and historical Markdown reports, obscuring the source/\allowbreak{}secondary-artifact boundary. & Root inspection of the V1.12 package found 36 Markdown reports, 20 PDFs, 20 TeX files, 19 aux files, 22 logs, 8 out files, 11 toc files, and additional latexmk artifacts at maxdepth 1. These files are not Lean sources and are reproducible or archival documentation surfaces. & Root clutter can make generated-secondary artifacts look like authoritative inputs, increases package size, and weakens the Single-YAML-Master /\allowbreak{} Lean-primary discipline. & Keep README.md in root;\newline archive historical Markdown reports;\newline remove generated PDF/\allowbreak{}TeX/\allowbreak{}LaTeX build products from root;\newline use scripts/\allowbreak{}generate\_full\_document.py for disposable full documentation and scripts/\allowbreak{}clean\_generated\_root\_artifacts.py for cleanup.\\
\hline
\addlinespace[0.24em]
SCF47 & \cellcolor{green!18}G & closed-by-V1.15-policy & 20.10 & V1.15 & 2026-05-10T12:20:00+02:00 & The V1.14 root policy was too narrow because Lean/\allowbreak{}Lake project-control files must live in the package/\allowbreak{}project root for reproducible local builds. & User supplied CNNA.lean, lakefile.lean, lake-manifest.json and lean-toolchain. The lakefile points the CNNA library to Repository/\allowbreak{}repo\_snapshot, while the manifest/\allowbreak{}toolchain pin the Lean/\allowbreak{}mathlib build environment. & If root cleanup deletes or omits these files, the package remains documentable but loses its direct local Lake build entry and reproducibility surface. & Amend DBR8/\allowbreak{}GNG27 by DBR10/\allowbreak{}GNG28: root is source/\allowbreak{}control oriented, not README-only. Keep package-root Lean/\allowbreak{}Lake build files;\newline continue excluding generated PDF/\allowbreak{}TeX/\allowbreak{}build artifacts and historical reports from active root.\\
\hline
\addlinespace[0.24em]
SCF48 & \cellcolor{green!18}G & closed-by-tooling-phase & 20.11 & V1.17 & 2026-05-10T12:58:00+02:00 & Root license/\allowbreak{}credit and duplicate-README surface needed explicit closure & V1.16 had README copies in both package root and Repository/\allowbreak{}, and no root LICENSE file. User required root-only README, root license file, and explicit editor/\allowbreak{}assistant credits also in generated PDFs. & Without closure, release packages could present ambiguous entry surfaces and generated PDFs could lose attribution/\allowbreak{}license context. & V1.17 keeps README.md and LICENSE only in package root, removes duplicate README files, updates PDF generators and metadata credits, and records GNG29/\allowbreak{}DBR11.\\
\hline
\addlinespace[0.24em]
SCF49 & \cellcolor{green!18}G & closed-by-tooling-phase & 20.12 & V1.18 & 2026-05-10T13:18:00+02:00 & Generated documentation sections needed stronger visual separation for review navigation & V1.17 full and context PDFs rendered consecutive top-level sections without a uniform automatic pagebreak invariant;\newline individual manual clearpages existed only around the table of contents or specific generator positions. & Large review packets become harder to navigate and section boundaries can be visually blurred, especially in context-filtered documents that preserve all major sections but show sparse rows. & V1.18 adds a generator-level section hook so every top-level full/\allowbreak{}context documentation section starts on a fresh page. The change is layout-only and does not alter Lean, YAML row semantics, context filtering, or proof status.\\
\hline
\addlinespace[0.24em]
SCF50 & \cellcolor{green!18}G & closed-by-tooling-phase & 20.13 & V1.19 & 2026-05-10T13:45:00+02:00 & White phase status was too permissive for unreleased Lean-derivation phases & V1.18 had many W phases with implementation\_scratchpad.phase\_type = derivation, so W still meant unchecked/\allowbreak{}future rather than non-Lean/\allowbreak{}tooling-only. User required unreleased Lean-derivation phases to default red until their release is proven from derived-only Lean objects. & Without this policy, a future Lean-content phase can look merely inactive/\allowbreak{}white instead of explicitly locked, weakening the derived-only partial-order discipline and hiding missing prerequisite-object evidence. & V1.19 converts all W derivation phases to R, redefines W for non-direct-Lean work, and patches generator validation to reject any future W derivation phase.\\
\hline
\addlinespace[0.24em]
SCF51 & \cellcolor{green!18}G & closed-by-tooling-phase & 20.14 & V1.20 & 2026-05-10T13:52:00+02:00 & Implementation module manifest could be misread as a second phase table and section/\allowbreak{}table purposes were underexplained & V1.19 still presented the module manifest in the main reading flow and several sections relied on terse table-local paragraphs. User reported confusion and requested clearer static explanations with dynamic table contents only. & Readers may confuse module/\allowbreak{}API change planning with phase-release logic, or miss where data comes from and why a table is shown. & V1.20 renames and moves the manifest to a technical appendix, clarifies its use boundary, and adds static source/\allowbreak{}purpose explanations around the full and context documentation sections.\\
\hline
\addlinespace[0.24em]
SCF52 & \cellcolor{yellow!25}Y & active-in-phase-1.1.1.1.5.1 & 1.1.1.1.5.1-1.1.1.1.5.5 & V1.21 & 2026-05-10T14:20:00+02:00 & Pre-code readiness was too broad: semantic consolidation alone did not prove that all current modules were neutralized, refactored and classified before new substantive Lean code. & V1.20 had active 1.1.1.1.5 and orange 1.1.1.2, while many required neutralization/\allowbreak{}refactor/\allowbreak{}object-intake phases remained white planning/\allowbreak{}refactor rows. The codebase still contains phase-semantic SM*/\allowbreak{}S6*/\allowbreak{}StrengtheningPhase names and 8860 declarations needing stable classification. & If 1.1.1.2 starts too early, blue-object correction may be built on unstable module names, incomplete IST-object intake, stale public/\allowbreak{}API surfaces or unclassified proof-stage modules. & V1.21 splits 1.1.1.1.5 into a five-step readiness chain and makes 1.1.1.2 depend on the final release certificate.\\
\hline
\addlinespace[0.24em]
SCF53 & \cellcolor{green!18}G & closed-as-planning-record & 20.15-20.20 & V1.22 & 2026-05-10T14:55:00+02:00 & Database/\allowbreak{}public-web vision needed explicit non-truth-source planning phases before any direct DB work. & SLTG-06 already described database-ready documentation, but provider-neutral SQL database public read-only mirror, snapshot provenance, read-only access policy/\allowbreak{}views and cnna-project.net frontend contract were not yet split into concrete planning phases. & Without explicit phases and guardrails, a future database could become an untracked second truth system, leak write authority, lose snapshot provenance or expose admin write credential secrets to browser/\allowbreak{}frontend contexts. & V1.22 adds 20.15-20.20 plus DBR15-DBR20 and GNG34 as a white non-blocking DB/\allowbreak{}tooling sidechain. Lean/\allowbreak{}YAML remain authoritative and active Lean/\allowbreak{}content cursor is unchanged.\\
\hline
\addlinespace[0.24em]
SCF54 & \cellcolor{green!18}G & closed-by-tooling-phase & 20.21 & V1.23 & 2026-05-10T15:20:00+02:00 & Generated PDFs needed a larger typography profile for review readability & V1.22 used a 10pt document class for full/\allowbreak{}context PDFs and many explicit scriptsize/\allowbreak{}footnotesize table fragments, which made large review packets dense even after section-pagebreak and explanation improvements. & Reviewers may miss status/\allowbreak{}provenance details in dense tables;\newline readability pressure could otherwise lead to manual PDF edits or inconsistent per-table hacks. & V1.23 applies a one-step font-size increase at generator level: documentclass 11pt, TeX size-command bump in tables, and graph label-size bump where generated by the tool.\\
\hline
\addlinespace[0.24em]
SCF56 & \cellcolor{green!18}G & closed-by-planning-phase & 1.4.9;\newline 13.2.1;\newline 13.6.1;\newline 18.5 & V1.25 & 2026-05-11T08:40:00+02:00 & Quaternion/\allowbreak{}geometric-algebra material needed reference-gate routing before it could influence CNNA planning. & Horn 0709.2238v1 contains useful Hamilton/\allowbreak{}Pauli/\allowbreak{}Clifford/\allowbreak{}duality/\allowbreak{}biquaternion target information, but using it directly would blur CNNA-derived-only provenance and sign/\allowbreak{}unit semantics. & Without explicit gates, external Pauli/\allowbreak{}Clifford/\allowbreak{}Dirac or biquaternion structures could silently become inputs, set orientation or scalar layers, or patch later Pillar-B physics from established theory. & V1.25 adds GA0-GA6 and comparison phases;\newline GNG37 forbids generator/\allowbreak{}handoff rights and keeps all material as target/\allowbreak{}scratchpad context only.\\
\hline
\addlinespace[0.24em]
SCF57 & \cellcolor{green!18}G & closed-by-V1.26-tooling-schema & 20.27 & V1.26 & 2026-05-11T09:25:00+02:00 & Phase/\allowbreak{}object status could be misread as context/\allowbreak{}type identity;\newline tool-only, database, generated-artifact, governance and reference-gate objects needed explicit separation from Lean objects. & Earlier generated phase/\allowbreak{}object tables showed status and links but did not render an explicit class field on all phase/\allowbreak{}object logic surfaces. & Category error: green generated or reference objects could be mistaken for Lean-derived objects;\newline red/\allowbreak{}yellow Lean phases could be conflated with non-Lean tooling phases. & Add Context class taxonomy and mandatory field, render in phase/\allowbreak{}object/\allowbreak{}proof/\allowbreak{}scratchpad/\allowbreak{}module-manifest tables, and enforce status/\allowbreak{}context orthogonality via GNG38.\\
\hline
\addlinespace[0.24em]
SCF58 & \cellcolor{green!18}G & closed-by-V1.29-three-layer-versioning & tooling/\allowbreak{}versioning-policy V1.29 & V1.29 & 2026-05-11T23:46:43+02:00 & Versioning still mixed Lean truth state, YAML planning state and tool-infrastructure state. & V1.28 separated tool/\allowbreak{}display version from Lean-data timestamp, but a planning-only YAML change and a generator-code change could still both appear mainly through the tool version, while phase/\allowbreak{}object/\allowbreak{}scratchpad/\allowbreak{}proof-dossier state had no independent planning-data version. & Reviewers could misread a plan-table update as a tool-code release, or a generator-code change as a Lean-code/\allowbreak{}data change. Snapshot provenance for YAML2SQL and the CNNA Explorer would also be ambiguous. & V1.29 introduced three independent metadata axes;\newline V1.30 tightens them into timestamp-complete pairs: lean\_data\_version/\allowbreak{}lean\_data\_timestamp, planning\_data\_version/\allowbreak{}planning\_data\_timestamp and tool\_infrastructure\_version/\allowbreak{}tool\_infrastructure\_timestamp. Legacy tool\_version, data\_version, data\_timestamp and tool\_change\_timestamp are retained only as compatibility aliases.\\
\hline
\addlinespace[0.24em]
SCF59 & \cellcolor{green!18}G & closed-by-V1.30-timestamp-complete-versioning & tooling/\allowbreak{}versioning-policy V1.30 & V1.30/\allowbreak{}P1.30 & 2026-05-11T23:53:41+02:00 & Version labels alone are insufficient for reliable history, SQL snapshots, Explorer provenance and generated artifact review. & After V1.29 the three axes existed, but downstream prose and some snapshot contracts could still emphasize the version value without always requiring the paired timestamp at every persistence boundary. & History becomes incomplete or ambiguous when a planning/\allowbreak{}tool/\allowbreak{}Lean version is stored, displayed, imported or compared without the exact timestamp of that axis. A later YAML2SQL snapshot or Explorer diff could then merge distinct states or fail to reconstruct change order. & V1.30 defines every semantic version axis as a mandatory version+timestamp pair, updates the snapshot/\allowbreak{}provenance wording, and makes the generators validate all six metadata fields: lean\_data\_version, lean\_data\_timestamp, planning\_data\_version, planning\_data\_timestamp, tool\_infrastructure\_version and tool\_infrastructure\_timestamp.\\
\hline
\addlinespace[0.24em]
SCF60 & \cellcolor{green!18}G & closed-by-V1.31-lean-toolchain-version-semantics & tooling/\allowbreak{}versioning-policy V1.31 & V1.31/\allowbreak{}P1.31 & 2026-05-11T23:59:59+02:00 & Lean-axis version and timestamp were still too similar and could be mistaken for duplicate snapshot identifiers. & V1.30 required lean\_data\_version + lean\_data\_timestamp, but lean\_data\_version still contained a timestamp-like value. This made the Lean-axis pair look redundant and obscured the distinction between Lean/\allowbreak{}toolchain/\allowbreak{}API level and Lean source/\allowbreak{}data snapshot time. & History, SQL snapshots and Explorer provenance could later fail to distinguish a Lean toolchain/\allowbreak{}API upgrade from a Lean source snapshot change. Reviewers could also assume both Lean-axis fields always move together. & V1.31 sets lean\_data\_version to the static Lean toolchain value lean-4.28.0, keeps lean\_data\_timestamp at the frozen Lean source snapshot time, and updates generator labels/\allowbreak{}prose so the Lean axis is rendered as Lean toolchain version + Lean data timestamp.\\
\hline
\addlinespace[0.24em]
SCF61 & \cellcolor{green!18}G & closed-planning/\allowbreak{}tooling-schema & 21.0 & P1.32 /\allowbreak{} V1.32 & 2026-05-12T00:14:05+02:00 & Object lifecycle/\allowbreak{}provenance gap: patched, renamed or deleted objects could previously be overwritten or disappear from current object views. & V1.32 adds object lifecycle fields, CEX10, non-destructive patch/\allowbreak{}rename/\allowbreak{}delete workflows, phase patch-link workflow, and generator rendering for lifecycle metadata. & Without lifecycle state and successor/\allowbreak{}current-dossier pointers, Explorer/\allowbreak{}SQL history could contradict prior provenance, hide old dossiers, or imply that a historical phase became invalid/\allowbreak{}outdated. & Treat object currency as lifecycle metadata orthogonal to object status. Objects can become superseded/\allowbreak{}renamed/\allowbreak{}deleted-retired;\newline phases remain historical events and link to patch/\allowbreak{}successor phases instead of becoming outdated.\\
\hline
\addlinespace[0.24em]
SCF62 & \cellcolor{green!18}G & closed-planning/\allowbreak{}tooling-schema & 21.0 /\allowbreak{} 1.1.1.1.5.1 & P1.33 /\allowbreak{} V1.33 & 2026-05-12T00:30:07+02:00 & Declaration consumption records did not guarantee total classification of the 8000+ Lean declarations by syntactic type, semantic role, triviality, documentation requirement, lifecycle and owner bucket. & Before P1.33, declaration\_consumption\_records.generated.tsv had about 8860 rows with consumption/\allowbreak{}module/\allowbreak{}public-promotion fields, while only registered CNNA objects had context/\allowbreak{}significance/\allowbreak{}dossier/\allowbreak{}lifecycle fields. This left a gap between source inventory and Explorer-ready declaration semantics. & The Explorer could display declarations as if fully sorted, or hide/\allowbreak{}overexpose declarations, while rfl/\allowbreak{}projection wrappers, legacy markers, public API surfaces, deleted/\allowbreak{}renamed declarations and fundamental definitional contracts were not uniquely classified. & P1.33 adds CEX11, total declaration classification taxonomy, workflows and guardrails. LEG37/\allowbreak{}CEX0 closure now requires exact declaration\_kind, semantic\_role, triviality\_tier, documentation policy, lifecycle\_state, public\_surface\_status and owner bucket for every declaration, plus manual override for heuristic labels.\\
\hline
\addlinespace[0.24em]
SCF63 & \cellcolor{green!18}G & closed-addressed-by-plan & 22;\newline 13.8-13.9;\newline 15.5-15.6;\newline 16.5-16.6;\newline 22.0-22.7 & P1.34/\allowbreak{}V1.34 & 2026-05-12T01:06:43+02:00 & Original Roadmap v0.0605 contained migration, chain-sheet, de-seeding, ParameterClosure, D/\allowbreak{}E and Feynman-audit obligations not yet fully represented in V1.33 planning. & Roadmap scan found missing or under-specified terms/\allowbreak{}structures: De-seeding, Kettenblatt, ParameterClosure, Nuclearity, DHR/\allowbreak{}DR, Feynman matter audit, local covariance/\allowbreak{}Hadamard/\allowbreak{}Wick, scattering/\allowbreak{}IR horizon, meta-audit markers and old-path rollback. & Without integration, future Explorer/\allowbreak{}SQL views and phase planning could look coherent while omitting critical legacy-roadmap gate obligations or semantic warnings. & Add roadmap-integration block 22, concrete future phases 13.8/\allowbreak{}13.9/\allowbreak{}15.5/\allowbreak{}15.6/\allowbreak{}16.5/\allowbreak{}16.6, objects B07/\allowbreak{}B08/\allowbreak{}PC00/\allowbreak{}D04/\allowbreak{}DE00/\allowbreak{}E04/\allowbreak{}E05/\allowbreak{}RDV0-RDV6, guardrails GNG46-GNG50 and generic external-source cross-check rendering.\\
\hline
\addlinespace[0.24em]
SCF64 & \cellcolor{green!18}G & closed-static-finding & 23.0 & P1.35 & 2026-05-12T01:19:51+02:00 & REALOQS legacy AQFT surface is substantial and must be treated as a first-class migration/\allowbreak{}reference source, not ignored. & Static scan found 270 Lean files, 1727 declarations, 88 PillarB files, 62 PillarB/\allowbreak{}AQFT files and 596 AQFT declarations in REALOQS v0.0605 source zip. & Without a ledger, advanced legacy AQFT work could be silently reintroduced or duplicated without provenance. & Add RLC0-RLC8 and block 23 to route legacy AQFT content through migration disposition and CNNA source maps.\\
\hline
\addlinespace[0.24em]
SCF65 & \cellcolor{red!20}R & open-blocker & 23.1-23.3 & P1.35 & 2026-05-12T01:19:51+02:00 & REALOQS AQFT code cannot be ported 1:1 into CNNA. & Legacy AQFT uses custom PillarA/\allowbreak{}BInterfaces, TreeOfCliques examples/\allowbreak{}configs, noncomputable sections and @[simp] patterns. The BoundaryMatrix derived layer is entirely noncomputable in the scan. & Direct reuse would violate derived-only, handoff filtering, operational-path and naming-neutrality constraints. & Require migration disposition, A->B dependency inversion, noncomputable isolation and neutral target modules before implementation.\\
\hline
\addlinespace[0.24em]
SCF66 & \cellcolor{yellow!25}Y & open-planned & 23.4-23.6 & P1.35 & 2026-05-12T01:19:51+02:00 & Legacy HK/\allowbreak{}GNS/\allowbreak{}KMS/\allowbreak{}quasi-local modules provide useful target gates but not CNNA closure. & Legacy code contains HK gates, StarAlgebra/\allowbreak{}State/\allowbreak{}GNS/\allowbreak{}KMS/\allowbreak{}quasi-local/\allowbreak{}split structures and derived boundary-matrix routes. & The names may suggest established AQFT closure even when CNNA source prerequisites are absent. & Record target gates and phase order under RLC4-RLC6;\newline force re-derivation from B local net and later C/\allowbreak{}D prerequisites.\\
\hline
\addlinespace[0.24em]
SCF67 & \cellcolor{red!20}R & open-blocker & 23.8 & P1.35 & 2026-05-12T01:19:51+02:00 & The legacy declarations need declaration-level crosswalk before Explorer visibility or migration claims. & Initial scan captures module/\allowbreak{}declaration counts but does not assign every legacy declaration to CNNA owner buckets, target gates, reject/\allowbreak{}defer buckets or lifecycle states. & Explorer could display legacy nodes as if they were CNNA-derived or current. & Make RLC8 dependent on CEX11 and block legacy overlay completeness until every visible legacy declaration has explicit non-CNNA epistemic status.\\
\hline
\addlinespace[0.24em]
SCF68 & \cellcolor{green!18}G & closed-by-P1.40-workflow-audit & 21.2.4 /\allowbreak{} workflow policy & P1.40 & 2026-05-12T08:18:00+02:00 & Workflow 7 could be read as setting a phase red whenever the assistant had not produced a green build, conflating absent build evidence with verified build failure. & P1.39 Implementation-from-last-green-baseline said otherwise set the attempted phase R, while Workflow 11 correctly says no build means artifact-generated /\allowbreak{} local Lean build required, not green. & The automated phase loop could create unnecessary repair phases or red-lock a phase merely because the current artifact package was planning-only or build evidence was missing. & P1.40 narrows red transitions to verified local/\allowbreak{}CI failure evidence. Missing build evidence stays evidence-pending/\allowbreak{}build-required and cannot close, red-lock or advance a phase by itself.\\
\hline
\addlinespace[0.24em]
SCF69 & \cellcolor{green!18}G & closed-by-P1.40-workflow-audit & 10 /\allowbreak{} 73 /\allowbreak{} workflow policy & P1.40 & 2026-05-12T08:18:00+02:00 & Post-green closeout still required generated tables and PDF compile, which could be misread as mandatory full-PDF regeneration after P1.39 introduced timestamp-bounded delta PDFs. & Workflow 10 predates the delta-PDF rule and did not distinguish bounded planning deltas from full release/\allowbreak{}audit checkpoints. & The documentation pipeline could reintroduce full-PDF churn and defeat the timestamp-bounded review discipline. & P1.40 rewrites closeout to require regenerated/\allowbreak{}validated tables and the applicable documentation artifact: delta PDF by default for bounded changes, full PDF only when explicitly requested or release-gated.\\
\hline
\addlinespace[0.24em]
SCF70 & \cellcolor{green!18}G & closed-by-P1.40-workflow-audit & 21.3 /\allowbreak{} 21.9 /\allowbreak{} CEX3 /\allowbreak{} CEX9 & P1.40 & 2026-05-12T08:18:00+02:00 & Red-build failure visibility existed as CEX16 but was not explicitly threaded through all downstream API and deployment dependencies. & P1.39 added CEX16 and 21.2.5, but 21.3/\allowbreak{}CEX3 and 21.9/\allowbreak{}CEX9 still named CEX12-CEX15 in several dependency fields. & Explorer API/\allowbreak{}deployment planning could expose green/\allowbreak{}current views while omitting red failure history or repair-phase publication obligations. & P1.40 adds CEX16 and 21.2.5 to 21.3/\allowbreak{}CEX3 and 21.9/\allowbreak{}CEX9 source chains, proof dossiers and scratchpads.\\
\hline
\addlinespace[0.24em]
\end{longtable}
\arrayrulecolor{black}
\normalsize
